
\section{Related Work}
\label{sec:related_work}

\subsection{Graph-Based Retrieval-Augmented Generation}

Recent advances in retrieval-augmented generation (RAG) have increasingly emphasized structural awareness to improve reasoning depth and contextual coherence. 
% A key direction involves constructing graph-structured indices from unstructured corpora to support richer navigation and inference. 
\citet{edge2024local} propose GraphRAG, which builds a knowledge graph (KG) from documents via LLM-based entity and relation extraction, then applies community detection to generate hierarchical summaries for global sensemaking. 
% This enables comprehensive answers to broad queries through multi-level summarization. 
\citet{guo2024lightrag} introduce LightRAG, which employs a dual-level retrieval system combining low-level fact retrieval and high-level semantic discovery using a compact KG, improving both efficiency and coverage. 
% Both methods rely on high-quality LLM-generated triples during indexing, making them sensitive to noise when using smaller models.
Further building on this idea, \citet{hipporag, gutiérrez2025HippoRAG2} present a non-parametric continual learning framework that uses Personalized PageRank over an open KG to enable associative, multi-hop reasoning. 
% It integrates passage-level context into the graph traversal process and leverages an LLM online to filter irrelevant triples, achieving strong performance across factual, sense-making, and associative memory tasks. Unlike standard RAG, HippoRAG 2 mimics human long-term memory by supporting dynamic, context-aware navigation through knowledge. 
% Our work draws inspiration from its use of graph diffusion and online LLM guidance but differs in introducing a multi-agent framework for subgraph refinement and unifying local and global retrieval via a heterogeneous graph.
Other structure-augmented RAG methods include RAPTOR~\citep{raptor}, \citet{chen2023walkingmemorymazecontext} enhance sense-making but often introduce noise through uncontrolled summarization or lack explicit support for multi-hop reasoning.
% Other structure-augmented RAG methods include RAPTOR \citep{raptor}, which recursively clusters and summarizes tree-like structures of text chunks to form a semantic hierarchy, enabling deeper context integration. Similarly, \citet{chen2023walkingmemorymazecontext} explore memory-augmented architectures that simulate associative recall. These approaches enhance sense-making but often introduce noise through uncontrolled summarization or lack explicit support for multi-hop reasoning.
Note that RL-based frameworks such as GraphRAG-R1~\citep{yu2025graphrag-r1} and Graph-R1~\citep{luo2025graph-r1} utilize existing Graph-based RAG methods as their retrieval components and train an end-to-end agentic framework. Our work, in contrast, proposes a complementary approach to improve the underlying retrieval paradigm itself. Consequently, ToG-3 could also serve as a plug-in component to enhance such RL frameworks.

\subsection{Knowledge Graphs in RAG and Hybrid Approaches}

The integration of structured knowledge into LLM reasoning has long been pursued to improve faithfulness and interpretability. 
Early KG-augmented RAG systems retrieve triples from static external knowledge bases such as Wikidata or Freebase to ground model outputs~\citep{sun2023tog}. However, these sources are often incomplete, outdated, or misaligned with domain-specific content. To overcome this, hybrid RAG frameworks~\citep{ma2024tog2} combine unstructured text and structured KGs to balance breadth and precision.
Chain-of-Knowledge (CoK)~\citep{li2024cok} retrieves from multiple structured sources including Wikipedia, Wikidata, and Wikitable to ground LLM responses. HybridRAG~\citep{sarmah2024hybridrag} fuses vector-based and KG-based retrievers, demonstrating superior reasoning performance compared to either modality alone. 
% However, most such systems adopt a \textit{loose-coupling} design, where retrieval from each source is independent and results are merged post-hoc, limiting synergistic interaction.
% In contrast, Think-on-Graph 2.0 (ToG-2) \citep{ma2024tog2} proposes a \textit{tight-coupling} paradigm: it uses a KG to guide document retrieval and leverages documents as contexts to refine graph traversal. This bidirectional interaction enables iterative, depth-first exploration of relevant knowledge, closely mimicking human reasoning. ToG-2 alternates between graph and context retrieval, allowing the LLM to progressively gather in-depth clues. While effective, it depends on pre-existing, high-quality KGs, restricting its applicability to domains where such resources are unavailable.

% Our work extends the tight-coupling principle of ToG-2 but removes the dependency on external KGs by dynamically constructing a heterogeneous graph from the document corpus. Furthermore, we enhance the iterative loop with explicit agent roles—Constructor, Reflector, and Response—enabling structured reflection and subgraph refinement.

\subsection{Iterative and Reflective Reasoning in LLMs}

Enabling LLMs to reason iteratively has been shown to improve accuracy and faithfulness. ITER-RETGEN~\citep{shao2023enhancingretrievalaugmentedlargelanguage} introduces an iterative loop that alternates between retrieval and generation, using generated hypotheses to guide further search. \citet{trivedi2023interleavingretrievalchainofthoughtreasoning} combine Chain-of-Thought (CoT) with retrieval, interleaving reasoning steps with evidence gathering, significantly improving performance on multi-hop QA. 
% However, these methods lack a structured memory substrate to ground the iteration, leading to redundancy and error propagation.
% To address this, several works integrate reasoning traces with retrieval indices. 
Self-RAG~\citep{asai2023selfraglearningretrievegenerate} equips LLMs with reflection tokens to decide when to retrieve and whether the output is hallucinated. ReAct~\citep{yao2023react} combines reasoning traces with external actions, enabling task decomposition and environment interaction. 
% These reflection mechanisms inspire our Reflector Agent, which explicitly evaluates the sufficiency of retrieved context and triggers sub-query generation and sub-graph refinement when needed.
Other efforts focus on continual learning for LLMs, where RAG serves as a non-parametric alternative to fine-tuning~\citep{shi2024continual}. Continual pretraining~\citep{lifelong} and instruction tuning~\citep{citb} can update model parameters but suffer from catastrophic forgetting~\citep{huang24mitigating}. Model editing methods~\citep{yao23editing} offer fine-grained updates but struggle with generalization. 
% In contrast, RAG externalizes knowledge, avoiding parameter updates altogether. ToG-3 advances this vision by organizing external memory as a dynamic, heterogeneous graph and enabling intelligent, agent-driven navigation—bringing LLMs closer to human-like, adaptive reasoning.

% In summary, ToG-3 integrates insights from graph-augmented RAG (GraphRAG, LightRAG), multi-hop reasoning (HippoRAG 2), tight-coupling design (ToG-2), and reflective generation (ITER-RETGEN, Self-RAG), while introducing a novel multi-agent optimization loop and a unified heterogeneous graph index. This allows it to achieve deep, adaptive, and deployable reasoning without relying on external KGs or high-resource indexing.

% Our proposed ToG-3 framework is inspired by these advancements. We inherit the iterative, reflective reasoning loop from the ToG series and apply it to a graph built directly from the corpus, like GraphRAG and LightRAG. Furthermore, our use of a multi-agent system and an implicit optimization loop to evolve the sub-query and refine the retrieved sub-graph is a novel approach that addresses the limitations of LLM-based graph construction, making our method more efficient and reliable, especially for smaller models. This allows ToG-3 to provide robust, high-quality reasoning that combines the strengths of iterative retrieval, corpus-based graph construction, and multi-agent systems to solve complex problems more effectively.